\clearpage
\fancyhead[L]{\small CERES. PALLAS. JUNO. COMET.}
\SourceLeaf{272}
\begin{center}
CERES, PALLAS, JUNO AND THE SECOND COMET OF 1805.
\end{center}

\noindent\rule{\linewidth}{0.4pt}
\begin{center}
Astronomisches Jahrbuch für 1809. Pp. 137–140. Berlin 1806.
\end{center}
\noindent\rule{\linewidth}{0.4pt}

\begin{center}
GAUSS TO BODE.
\end{center}
\hfill Brunswick, 14 March 1806.

\par\indent I am pleased to pass on to Your --- a few observations that I made last winter, which I hope you will not find unwelcome. There are few of them, but the bad weather that all my astronomical correspondents complain of has prevailed here too throughout the winter. Besides, with the scant means at my disposal, I can pursue practical work here only as a secondary occupation. First come my observations of the three new planets.

\Needspace{13\baselineskip}
\begin{center}\emph{Ceres.}\end{center}
\begin{center}
\begingroup\fontsize{9}{11}\selectfont\setlength{\tabcolsep}{2pt}\renewcommand{\arraystretch}{1.1}
\begin{tabular}{lr|r|r}
\multicolumn{2}{c|}{Mean time} & R.A. & Decl.\\ \hline
\textsuperscript{1}805. 24 October & $11^{\mathrm u}\ 22^{\mathrm m}\ 32^{\mathrm s}$ & $109^\circ\ 28'\ 29''$ & $23^\circ\ 27'\ 0''$ North.\\
28. & $11\ 31\ 22$ & $110\ 7\ 7$ & $23\ 36\ 58$\\
14 November & $9\ 22\ 30$ & $111\ 44\ 56$ & $24\ 32\ 1$\\
18. & $9\ 0\ 37$ & $111\ 50\ 48$ & $24\ 48\ 47$\\
\end{tabular}\endgroup
\end{center}

\Needspace{11\baselineskip}
\begin{center}\emph{Pallas.}\end{center}
\begin{center}
\begingroup\fontsize{9}{11}\selectfont\setlength{\tabcolsep}{2pt}\renewcommand{\arraystretch}{1.1}
\begin{tabular}{lr|r|r}
1806. 14 February & $8^{\mathrm u}\ 11^{\mathrm m}\ 16^{\mathrm s}$ & $70^\circ\ 16'\ 31''$ & $19^\circ\ 59'\ 13''$ South.\\
16. & $7\ 32\ 28$ & $70\ 42\ 39$ & $19\ 20\ 44$\\
17. & $8\ 52\ 38$ & $70\ 56\ 44$ & $19\ 1\ 8$\\
20. & $7\ 49\ 35$ & $71\ 39\ 2$ & $18\ 5\ 0$\\
\end{tabular}\endgroup
\end{center}

\Needspace{8\baselineskip}
\begin{center}\emph{Juno.}\end{center}
\begin{center}
\begingroup\fontsize{9}{11}\selectfont\setlength{\tabcolsep}{2pt}\renewcommand{\arraystretch}{1.1}
\begin{tabular}{lr|r|r}
1806. 17 February & $9^{\mathrm u}\ 42^{\mathrm m}\ 0^{\mathrm s}$ & $173^\circ\ 46'\ 45''$ & $0^\circ\ 28'\ 32''$ North.\\
20. & $11\ 24\ 15$ & $173\ 15\ 53$ & $0\ 54\ 30$\\
\end{tabular}\endgroup
\end{center}

\par\indent According to these observations, for Ceres my ephemeris gave a right ascension too large by 4$\frac{1}{2}'$ in Oct. 1805 and by 5$\frac{1}{2}'$ in Nov., and a declination too small by 1$'$. For Pallas it gave a right ascension too large by 27$\frac{1}{2}'$ and a declination too large by 5$'$ (that is, too far north). For Juno it gave a right ascension too large by 1$\frac{1}{2}'$ and a declination too small by 1$'$ (that is, too far south). [These observations of Ceres are also in Monatl. Corr., vol. XIII, \href{https://www.digitale-sammlungen.de/en/view/bsb10538605?page=197}{p. 189}.]



\par\indent The considerable discrepancy for Pallas is undoubtedly due to the perturbations that are now beginning to show. For Juno, on the other hand, the agreement is remarkably good, beyond all expectation, and proves the accuracy of my latest elements and of my observation on 20 Feb. 1805, on which those elements were based. So far I know of no one else who has observed or seen Pallas or Juno this year. Pallas is very bright, as bright as a star of magnitude 7., or of 7..8\textsuperscript{th} magnitude. It must be easy to observe with a transit instrument and a quadrant even under full illumination. Juno, of course, is faint, about 11\textsuperscript{th} magnitude when I observed it, and so it can hardly be observed with an illuminated telescope. Ceres is as clearly visible as ever. The only reason I have been unable to continue observing it is that it rises too high for my observing site. Prof. \textsc{Harding} has often observed it successfully with the mural quadrant. I will be pleased to send you my results and Ceres’s path for the next apparition.

\SourceLeaf{273}

\par\indent I observed last year’s second comet only on 8 December, when I saw it twice. It was clearly visible even to the naked eye, about as bright as a star of the 3rd or 4th magnitude. Here are my observations.
\begin{center}
\begingroup\fontsize{9}{11}\selectfont\setlength{\tabcolsep}{2pt}\renewcommand{\arraystretch}{1.1}
\begin{tabular}{lr|r|r}
\multicolumn{2}{c|}{Mean time} & Right asc. & Decl.\\ \hline
1805. 8 December & $6^{\mathrm u}\ 55^{\mathrm m}\ 17^{\mathrm s}$ & $353^\circ\ 7'\ 40''$ & $23^\circ\ 36'\ 24''$ South.\\
& $7\ 58\ 3$ & $352\ 57\ 57$ & $24\ 6$\\
\end{tabular}\endgroup
\end{center}

\par\indent The parabolic elements I calculated from this observation, one by \textsc{Olbers} on 2 Dec., and the one by \textsc{Bouvard} on 16 Nov. published in the Moniteur, appear in the January issue of the M. C. [\hyperlink{comet1805-2.1}{p. 266} of this volume]. Since then, through the kindness of Mr. \textsc{von Zach}, I have also received the complete series of \textsc{Thulis}’s original observations, and the newspapers have reported a meridian observation made at Greenwich on 8 Dec. With these additional observations I obtained the following improved elements.

\Needspace{11\baselineskip}
\begin{center}
\begingroup\fontsize{8}{10}\selectfont\setlength{\tabcolsep}{2pt}\renewcommand{\arraystretch}{1.1}
\begin{tabular}{lr}
Time of perihelion 1805. 31 Dec. . . . . . & $6^{\mathrm u}\ 54^{\mathrm m}\ 35^{\mathrm s}$ Seeberg mean time\\
Longitude of the point nearest the Sun . . . . . . . . . & $109^\circ\ 21'\ 50''5$\\
Longitude of the ascending node . . . . . . & $250\ 33\ 34.9$\\
Inclination of the orbit . . . . . . . . . . . & $16\ 30\ 31.9$\\
Logarithm of the least distance . . . . & $9.9503300$\\
Motion . . . . . . . . . . . . . . & prograde.\\
\end{tabular}\endgroup
\end{center}

\par\indent Here is the comparison of all the observations with these elements. I did receive a few more from \textsc{Thulis}, but I could not reduce them because the stars used for comparison do not appear in the star catalogues or in the Hist. Cél. [The times for \textsc{Gauss}, \textsc{Olbers} and \textsc{Thulis} are those at Brunswick, Bremen and Marseille, respectively.]
\begin{center}
\begingroup\fontsize{8}{10}\selectfont\setlength{\tabcolsep}{2pt}\renewcommand{\arraystretch}{1.1}
\begin{tabular}{r|r|r|r|r|r|l}
1805 & Mean time & \shortstack{Calculated right\\ascension} & \shortstack{Calculated\\declination} & \shortstack{Diff. in R.A.\\R---B} & \shortstack{Diff. in decl.\\R---B} & Observer\\ \hline
15 Nov. & $10^{\mathrm u}\ 28^{\mathrm m}\ 13^{\mathrm s}$ & $14^\circ\ 33'\ 0''$ & $39^\circ\ 18'\ 41''$ N. & $-2'\ 30''$ & $+\phantom{1'}\ 26''$ & Thulis\\
16 & $7\ 17\ 54$ & $14\ 12\ 7$ & $38\ 58\ 53$ & $-1\ 58$ & $+\phantom{1}\ 46$ & Thulis\\
19 & $7\ 11\ 17$ & $12\ 55\ 12$ & $37\ 34\ 27$ & $-\phantom{1}\ 4$ & $+1'\ 7$ & Thulis\\
20 & $7\ 3\ 15$ & $12\ 28\ 23$ & $36\ 59\ 41$ & $-\phantom{1}\ 23$ & $-\phantom{1}\ 32$ & Thulis\\
21 & $9\ 36\ 1$ & $11\ 57\ 5$ & $36\ 16\ 16$ & $+1\ 13$ & $-\phantom{1}\ 20$ & Thulis\\
29 & $6\ 39\ 44$ & $7\ 24\ 39$ & $26\ 47\ 49$ & $+1\ 32$ & $+\phantom{1}\ 1$ & Thulis\\
2 Dec. & $5\ 34\ 9$ & $4\ 48\ 43$ & $18\ 58\ 18$ & $+1\ 35$ & $-\phantom{1}\ 52$ & Olbers\\
2 & $9\ 39\ 0$ & $4\ 36\ 3$ & $18\ 20\ 28$ & $+1\ 47$ & $-\phantom{1}\ 3$ & Thulis\\
3 & $5\ 48\ 39$ & $3\ 39\ 51$ & $15\ 6\ 38$ & $+\phantom{1}\ 45$ & $-\phantom{1}\ 19$ & Olbers\\
3 & $6\ 30\ 37$ & $3\ 36\ 52$ & $14\ 55\ 11$ & $+1\ 10$ & $-\phantom{1}\ 26$ & Thulis\\
4 & $6\ 31\ 5$ & $2\ 16\ 36$ & $10\ 6\ 29$ N. & $+2\ 12$ & $-1\ 12$ & Thulis\\
6 & $7\ 16\ 16$ & $358\ 39\ 22$ & $3\ 39\ 32$ S. & $+\phantom{1}\ 22$ & $+\phantom{1}\ 48$ & Thulis\\
7 & $7\ 35\ 21$ & $356\ 9\ 21$ & $13\ 4\ 56$ & $\phantom{+1}\ 0$ & $+\phantom{1}\ 42$ & Thulis\\
8 & $5\ 27\ 22$ & $353\ 20\ 30$ & $22\ 56\ 24$ & $-1\ 10$ & $+2\ 44$ & Olbers\\
8 & $6\ 46\ 23$ & $353\ 8\ 33$ & $23\ 34\ 40$ & $-\phantom{1}\ 42$ & $+2\ 18$ & Olbers\\
8 & $6\ 55\ 17$ & $353\ 7\ 44$ & $23\ 35\ 57$ & $+\phantom{1}\ 4$ & $-\phantom{1}\ 47$ & Gauss\\
8 & & $353\ 6\ 21$ & $23\ 40\ 59$ & $-\phantom{1}\ 33$ & $-\phantom{1}\ 15$ & Maskelyne\\
8 & $6\ 59\ 39$ & $353\ 3\ 54$ & $23\ 47\ 34$ & $+\phantom{1}\ 9$ & $-\phantom{1}\ 33$ & Thulis\\
8 & $7\ 58\ 3$ & $352\ 57\ 34$ & $24\ 6\ 39$ & $-\phantom{1}\ 23$ & & Gauss\\
\end{tabular}\endgroup
\end{center}

\par\indent The elements of this comet are indeed very close to those of the comet of 1772, but the difference, especially in the least distance, is still too great. My own careful exami%
\SourceLeaf{274}%
nation of the observations of the comet of 1772 also showed me that the two sets of elements cannot be brought substantially closer together without spoiling their agreement with the observations. Nor did the comet of 1772 come close enough to any major planet for that planet’s perturbation to cause a large change in its elements. [According to \textsc{von Zach}, \textsc{Gauss} communicated these remarks on 21 Feb. 1806 (Monatl. Corr., vol. XIII, \href{https://www.digitale-sammlungen.de/en/view/bsb10538605?page=320}{pp. 310–313}, March 1806) and at the same time gave some preliminary indications of work whose results are given more fully in the letter of 20 May 1806 printed below on \hyperlink{native-275-lower}{p. 270}.]

\par\indent I am now working on my method for determining planetary orbits from observations, and I hope before long to be able to think about making it public. The many improvements I have had the opportunity to introduce since 1801 have left its various parts bearing almost no resemblance to their original form. So I do not regret having taken my time over its publication.

